\documentclass[twocolumn]{iopjournal}

\usepackage{amsmath,amssymb}
\usepackage{booktabs}
\usepackage{array}
\usepackage{longtable}
\usepackage{hyperref}
\usepackage{xcolor}
\usepackage[small]{caption}

\fancyhfoffset[L]{0mm}
\setlength{\marginparwidth}{0pt}
\setlength{\marginparsep}{0pt}

% ==========================================
% ZENODO PREPRINT HEADERS (ACTIVE)
% ==========================================
\renewcommand{\articletype}[1]{%
  {\vspace*{-8mm}\noindent \Large \sf Preprint}\par
  \vspace*{6mm}\noindent{\scriptsize \sf{\bfseries \MakeUppercase{#1}}}}

\fancyhead[L]{}
\fancyhead[R]{}

% ==========================================
% CQG SUBMISSION HEADERS (COMMENTED OUT)
% ==========================================
% \fancyhead[L]{{\small \sf IOP Publishing}\hspace{5mm} {\it Class. Quantum Grav.}}
% \fancyhead[R]{S Nelson}

\newcommand{\FZERO}{F_{\mathrm{ZERO}}}
\newcommand{\FPLUS}{F_{\mathrm{PLUS}}}
\newcommand{\FMINUS}{F_{\mathrm{MINUS}}}
\newcommand{\BEA}{\mathrm{BE}/A}

\begin{document}

\articletype{Letter}

\title{Nuclear shell structure and binding energies from tetrahedral packing at $K\!=\!4$}

\author{Stephen Nelson}

\affil{Field Dynamics Research}

\begin{abstract}
A single algebraic formula, derived from the tetrad geometry of
general relativity, reproduces nuclear binding energies across all
83 stable elements from hydrogen to bismuth with zero adjusted
parameters.
The binding energy per nucleon is
$\BEA = C\,(f - \FZERO)$,
where $f(K,s) = 1 + Ks + K^2$ is the coupling polynomial of the
stellated octahedron evaluated at nuclear compression,
$\FZERO = K^2 + 1 = 17$ is the Field Origin capacity, and the
energy constant $C = m_e c^2\,\FZERO / K = 2.172$~MeV is fixed by
the electron mass and the $K = 4$ topology.
Across 82 isotopes from helium to bismuth the formula achieves
$R^2 = 0.99$ on total binding energy, with mean error $3.9\%$
and $98\%$ of elements within $10\%$.
A packing algebra based on tetrahedral close-packing
with Sierpi\'nski gate exclusion produces all seven nuclear magic
numbers $\{2, 8, 20, 28, 50, 82, 126\}$ exactly, with $K = 4$ as the
unique solution among all integers $2 \leq K \leq 10$.
The nesting ceiling $K^K = 256$ predicts that no shell beyond the
eighth can close, consistent with bismuth-209 as the heaviest
stable isotope.
The same polynomial, evaluated at galactic scales with a single-face
coupling path, produces galaxy rotation curves without dark matter.
\end{abstract}


% ================================================================
\section{Framework and coupling polynomial}
% ================================================================

\subsection*{The tetrad is the tetrahedron}

The tetrad formulation of general relativity,
inherited from Einstein~(1928), Cartan~(1922), and
Weyl~(1929), describes the local gravitational frame
as a set of four independent basis vectors at each point
in spacetime.
Algebraically, it is a transformation matrix.
Geometrically, four non-coplanar vectors from a common origin
define a tetrahedron: the simplest solid enclosing volume
in $D = 3$ spatial dimensions, with $K = D + 1 = 4$ faces.
The tetrad is the algebra; the tetrahedron is the geometry.
They are two descriptions of the same local frame.
Regge~\cite{regge1961} made the identification explicit:
in his discretisation of general relativity, a 3-simplex
(a tetrahedron) is the fundamental building block of space,
and all curvature concentrates at the vertices where
tetrahedra meet.

This identification is not limited to the single-metric case.
Hassan and Rosen~\cite{hassanrosen2012} proved that ghost-free
bimetric gravity requires two independent tetrads,
$e^{a}_{\mu}$ and $f^{a}_{\mu}$, coupled through an interaction
potential. Since each tetrad is geometrically a tetrahedron,
the theory requires two tetrahedra at every point in spacetime.
Three further results reach the same conclusion:
vacuum stability requires opposite orientations to cancel
the Boulware--Deser ghost~\cite{boulwaredeser1972};
the $SL(2,\mathbb{C})$ double cover of the Lorentz group
assigns one chirality per chamber~\cite{weinberg1995};
and Ashtekar's self-dual/anti-self-dual
decomposition~\cite{ashtekar1986} assigns one chirality per
chamber.
Four independent arguments. Two tetrahedra.

Gravity Field Dynamics (GFD)~\cite{nelson2025a} begins here.
It takes the algebraic requirement for two interacting tetrads
and draws the geometry: what does the interior field structure
of a gravitational field look like when both tetrahedra are
present?

\subsection*{Two chambers, not one}

Consider how gravitational fields couple to their neighbours.
In any close-packed field lattice, each field region borders
twelve others, the face-centred cubic (FCC) coordination
number, established by Kepler in 1611 and proved optimal by
Hales in 1998. These twelve neighbours press inward on the
central field through shared boundary planes, the same way
twelve soap bubbles pressed around a central one flatten its
walls into planar facets.
The twelve boundary planes are not randomly oriented.
In FCC packing, they sort into exactly four groups of three,
each group sharing a common normal direction.
Four groups, four normals, four faces: a tetrahedron emerges
from the packing geometry itself.

But packing is symmetric. The twelve neighbours do not push
from one side only; they push from all sides. The four
face-normals of the inward-pointing tetrahedron have four
partners pointing outward in the opposite directions.
These outward normals define a second tetrahedron, inverted
with respect to the first.
The two tetrahedra interlock with opposite orientation,
each vertex of one sitting at the centre of a face of the other.
Their compound is the \emph{stellated octahedron}:
8~vertices, 12~edges, 8~triangular faces.

The interior of this compound has three distinct regions.
The primary tetrahedral chamber (TetA) is the outward-facing
field structure, the sector that couples to neighbours and
carries the massless gravitational mode.
The inverted chamber (TetB) is the bounded interior, the
sector that can hold confined field content and carries the
massive mode.
At the geometric centre, all four body diagonals cross at a
single point.
This is the \emph{Field Origin}: the coupling boundary between
the two chambers, the unique location where both tetrads are
evaluated at the same spacetime event.
It is the interaction vertex of bimetric
gravity~\cite{hassanrosen2012}, given geometric location.

On each face of TetA, the opposing vertex of TetB
penetrates the surface and subdivides it into $3 + 1 = 4$
sub-regions: three spatial corners and one central core.
Across all four faces, this produces $4 \times 4 = K^2 = 16$
independent coupling channels.
These channels are the $K^2$ components of the frame
transition matrix in bimetric gravity, now visible as
geometric features of the stellated octahedron.

\subsection*{The coupling polynomial}

The tetrad equation contains three structural types. A nondegeneracy condition ($\det(e) \neq 0$): one
binary constraint ensuring the field closes.
A frame ($e^{a}_{\mu}$): $K = 4$ independent directions spanning
the tangent space. A connection ($\eta_{ab}$): a $K \times K$
matrix with $K^2 = 16$ slots encoding how every direction
couples to every other.
One constraint, four legs, sixteen components.
The coupling polynomial counts these structural levels:
\begin{equation}
  f(K, s) = 1 + Ks + K^2\,,
  \label{eq:poly}
\end{equation}
where $s \in [-1, +1]$ parameterizes the position between the
two chambers. In the three structural regions of the stellated
octahedron, the polynomial evaluates to three integers:
\begin{align}
  \FPLUS  &= f(K, {+}1) = 21 \;\;(\text{TetA})\,, \nonumber\\
  \FZERO  &= f(K, 0)    = 17 \;\;(\text{Field Origin})\,, \nonumber\\
  \FMINUS &= f(K, {-}1) = 13 \;\;(\text{TetB})\,.
  \label{eq:landmarks}
\end{align}
These are arithmetic-locked: $(21 + 13)/2 = 17$, equally spaced
by $K = 4$, with no freedom to adjust.
The integer $K = 4$ is itself selected uniquely by
$2(D{+}1) = 2^D$, whose only positive integer solution is
$D = 3$~\cite{nelson2025a}.

The constitutive law $\mu(x) = x/(x{+}1)$ is not postulated
as in MOND phenomenology but derived as the M\"obius map of this
polynomial~\cite{nelson2025a}, and its antiderivative yields the
Bekenstein--Milgrom energy functional~\cite{bekenstein1984} that
produces galaxy rotation curves.
The Hassan--Rosen interaction potential, whose five free
parameters $\{\beta_0,\ldots,\beta_4\}$ are conventionally
fitted to data, is locked by the dual tetrahedral geometry to
$\{1, 0, K, 0, 1\}$~\cite{nelson2025a}, reducing five constants
to zero.
The framework is covariant-complete: the polynomial, the
constitutive law, the energy functional, and the bimetric
parameters all follow from $K = 4$ and the stellated octahedron.

\subsection*{Nuclear scales}

In this letter, we apply the coupling polynomial at nuclear
scales. TetA is the unbound gravitational field: spacetime
itself, the ocean in which matter sits. TetB is the bounded
interior: a cavity that can hold confined field content.
When stellar compression packs gravity fields into the TetB
cavity of a host field, the coupling path traverses the
complete Field Origin, crossing $N = 2$ faces.
The $N = 2$ constitutive law is
\begin{equation}
  x(1 - \mu)^2 = \mu\,,
  \label{eq:n2}
\end{equation}
with the physical root
\begin{equation}
  \mu_2(x) = \frac{(2x + 1) - \sqrt{4x + 1}}{2x}\,.
  \label{eq:mu2}
\end{equation}
The nuclear conformal map is $s_2(x) = 2\mu_2(x) - 1$,
and the compression ratio is $x = A^2$, where $A$ is the number
of gravity fields packed in the host cavity (the mass number).

Substituting $s = s_2(A^2)$ into $f(K, s)$ and subtracting the
Field Origin capacity gives the binding energy per field:
\begin{equation}
  \boxed{\;
    \frac{\mathrm{BE}}{A} = C \cdot \bigl[\,f\!\left(K,\; s_2(A^2)\right)
    - \FZERO\,\bigr]
  \;}
  \label{eq:binding}
\end{equation}
The coupling polynomial evaluated at the nuclear conformal map,
minus its value at the Field Origin, is the binding energy.
The range of $f$ from $\FMINUS = 13$ to $\FPLUS = 21$ maps to a
binding energy range of $C \times ({-}4)$ to $C \times (+4)$~MeV,
with zero binding at $f = \FZERO = 17$.


% ================================================================
\section{The energy constant}
% ================================================================

The energy constant $C$ sets the scale of nuclear binding in MeV.
We derive it from the electron mass and the topology.

The electron mass $m_e$ is the unique mass for which the gravity
envelope radius coincides with the classical electron radius,
derived from $K$, $G$, $\alpha$, $\hbar$, $c$, and
$a_0$~\cite{nelson2025a}.
The binding energy ceiling is the maximum energy the Field Origin
can channel into packing.
This is the electron rest energy (the vortex inversion energy
at the envelope boundary) multiplied by the gate capacity $\FZERO$,
distributed across $K$ active coupling channels:
\begin{equation}
  C = \frac{m_e c^2 \cdot \FZERO}{K}
    = \frac{0.511 \times 17}{4}
    = 2.172~\text{MeV}\,.
  \label{eq:C}
\end{equation}
The empirical value obtained by fitting to the 82-isotope mean
is $C_{\mathrm{fit}} = 2.16$~MeV, a discrepancy of $0.54\%$.
The asymptotic binding ceiling is
\begin{align}
  C(\FPLUS {-} \FZERO) &= C \times 4 \nonumber\\
  &= m_e c^2 \!\cdot\! \FZERO \nonumber\\
  &= 8.687~\text{MeV}\,,
\end{align}
consistent with the observed peak at Fe-56 ($\BEA = 8.790$~MeV),
whose individual element-specific binding exceeds the geometric
mean by approximately $1.2\%$~\cite{nelson2025b}.


% ================================================================
\section{Binding energy validation}
% ================================================================

Equation~\eqref{eq:binding} with the derived $C$ from
Eq.~\eqref{eq:C} is applied to 82 stable isotopes
(He-4 through Bi-209). No parameters are adjusted.

The results are tabulated in Table~\ref{tab:full}.
Summary statistics:
\begin{center}
\small
\begin{tabular}{lc}
\toprule
Quantity & Value \\
\midrule
Elements tested & 82 \\
$R^2$ (total BE) & 0.991 \\
Mean $|\mathrm{error}|$ (BE/A) & $3.9\%$ \\
Within $5\%$ & 65 ($79\%$) \\
Within $10\%$ & 80 ($98\%$) \\
Beyond $15\%$ & 1 (He-4) \\
Fitted parameters & 0 \\
\bottomrule
\end{tabular}
\end{center}

The single outlier is He-4, whose binding energy per field
($7.074$~MeV observed vs $4.853$~MeV predicted) is systematically
underpredicted. He-4 is the first complete $K = 4$ closure
($A = K$) and receives a geometric resonance bonus that the
constant-$C$ formula does not capture.
Excluding He-4, the mean error across 81 isotopes is $3.5\%$.

The heavy-element end (Pb, Bi, Tl, Hg) shows a systematic
overprediction of $6$--$10\%$, attributable to Coulomb repulsion
among vortex-active (proton) fields, which grows as $Z^2$ and is
not captured by the $A$-only formula.
The Bethe-Weizs\"acker mass formula achieves comparable coverage
with five fitted parameters.
This formula uses zero.


% ================================================================
\section{The packing algebra}
% ================================================================

The binding energy formula determines \emph{how tightly} fields
are held. The packing algebra determines \emph{how they arrange}:
which configurations are stable, where the magic numbers fall,
and why the shell structure has the form it does.

Gravity fields packed inside the tetrahedral cavity (TetB) arrange
in three-dimensional tetrahedral close-packing, constrained by the
$K = 4$ channel structure of the host. The packing is recursive:
once $K = 4$ fields fill the host channels, additional fields
pack inside the existing residents, each of which is itself a
complete stellated octahedron.
At each nesting level, the Field Origin of each host is reserved
for structural coupling to its parent (the gate). This gate is
not available for packing.
The gate exclusion follows the Sierpi\'nski construction:
at each triforce level, the central sub-triangle is removed,
leaving three outer sub-triangles for further recursion.

Three functions encode the complete shell structure.
The tetrahedral capacity is
\begin{equation}
  T(n) = \frac{n(n{+}1)(n{+}2)}{6}\,,
  \label{eq:T}
\end{equation}
the gate exclusion is
\begin{equation}
  G(n) = \begin{cases}
           0 & n < K \\[2pt]
           (n{-}1)\,n/2 & n \geq K
         \end{cases}\,,
  \label{eq:G}
\end{equation}
and the available positions are
\begin{equation}
  P(n) = T(n) - G(n)\,.
  \label{eq:P}
\end{equation}
The magic number at shell $n$ is
\begin{equation}
  \boxed{\; M(n) = 2\,P(n) \;}
  \label{eq:magic}
\end{equation}
where the factor of~2 accounts for the two time-face orientations
(spin degeneracy). The packing is subject to the nesting ceiling
\begin{equation}
  M(n) \leq K^K = 256\,.
  \label{eq:ceiling}
\end{equation}

Evaluating Eqs.~\eqref{eq:T}--\eqref{eq:magic} at $K = 4$:

\begin{center}
\footnotesize
\begin{tabular}{cccccc}
\toprule
$n$ & $T$ & $G$ & $P$ & $M$ & $S$ \\
\midrule
1 &   1 &  0 &   1 &    2 &  2 \\
2 &   4 &  0 &   4 &    8 &  6 \\
3 &  10 &  0 &  10 &   20 & 12 \\
4 &  20 &  6 &  14 &   28 &  8 \\
5 &  35 & 10 &  25 &   50 & 22 \\
6 &  56 & 15 &  41 &   82 & 32 \\
7 &  84 & 21 &  63 &  126 & 44 \\
8 & 120 & 28 &  92 &  184 & 58 \\
9 & 165 & 36 & 129 & \textbf{258} & 74 \\
\bottomrule
\end{tabular}
\end{center}
Gate exclusion ($G > 0$) begins at $n = K = 4$.
Shell~9 requires $M(9) = 258 > K^K = 256$ and cannot close.

The column $S(n) = M(n) - M(n{-}1)$ gives the shell size
(nucleons per shell). Shell~3 has $S = 12$ and shell~4 has
$S = 8$. This inversion ($S_3 > S_4$) is the spin-orbit effect:
in standard nuclear physics it requires introducing the spin-orbit
force with a fitted coupling constant.
Here it is a theorem.
Below $K$, shell sizes follow $n(n{+}1)$, giving $3 \times 4 = 12$
at $n = 3$. The transition shell at $n = K$ has size $2K = 8$.
Since $D \times K = 12 > 2K = 8$ (because $D = 3 > 2$), the
inversion is geometrically forced.

\paragraph{Uniqueness.}
Testing all integers $2 \leq K \leq 10$, only $K = 4$ produces
the observed magic number sequence
$\{2, 8, 20, 28, 50, 82, 126\}$.
$K = 3$ misses 28.
$K = 5$ misses 20.
The data selects the stellated octahedron uniquely.


% ================================================================
\section{The stability ceiling}
% ================================================================

The recursive nesting has maximum depth $K = 4$, corresponding to
the $3 + 1$ degrees of freedom of spacetime: three spatial
dimensions and one temporal. Each nesting level consumes one
degree of freedom.
The total addressable state space is $K^K = 4^4 = 256$.

Shell~8 closes at $M(8) = 184 < 256$, safely within the ceiling.
This corresponds to the predicted island of stability near
$Z \approx 114$, $N = 126$--$184$.
Shell~9 would require $M(9) = 258 > 256$, exceeding the ceiling
by exactly~2 (the number of time faces).
Shell~9 cannot close.
Beyond this point, nuclear configurations are progressively
unstable.

Bismuth-209, the heaviest observationally stable isotope, uses
$209/256 = 81.6\%$ of the address space.
Fermium-258 ($A = 258 = M(9)$) decays by spontaneous fission
with a half-life of $370\;\mu$s, consistent with complete
exhaustion of the nesting hierarchy.


% ================================================================
\section{Galactic and nuclear scales from one polynomial}
% ================================================================

The coupling polynomial $f(K, s)$ governs both scales of the
framework through a single structural parameter: $N$, the number
of tetrahedral faces on the coupling path.

At galactic scales ($N = 1$, one face crossing), the conformal map
is $s_1(x) = (x{-}1)/(x{+}1)$ and the constitutive law
$\mu_1(x) = x/(x{+}1)$ is the M\"obius map of the polynomial.
Its antiderivative
\begin{equation}
  F(y) = y - 2\sqrt{y} + 2\ln(1 + \sqrt{y})
\end{equation}
is the AQUAL energy functional~\cite{bekenstein1984}, from which
galaxy rotation curves follow without dark matter.
At nuclear scales ($N = 2$, full Field Origin traversal), the
conformal map shifts to $s_2(A^2) = 2\mu_2(A^2) - 1$, and the
polynomial evaluates to
\begin{align}
  f\bigl(K, s_2(A^2)\bigr) &= \FMINUS + 2K\,\mu_2(A^2) \nonumber\\
  &= 13 + 8\,\mu_2(A^2)\,.
\end{align}
The binding energy is $C$ times the excess above $\FZERO = 17$.

The two scales share the polynomial, the three landmark values
$(13, 17, 21)$, and the integer $K = 4$.
They differ only in $N$.
One face crossing for galaxies. Two for nuclei.
The same stellated octahedron at both scales.


% ================================================================
\section{Discussion}
% ================================================================

\paragraph{What is established.}
The binding energy formula Eq.~\eqref{eq:binding} with
derived $C$ (Eq.~\ref{eq:C}) reproduces 82 isotope binding
energies at $R^2 = 0.99$ with zero fitted parameters.
The packing algebra Eqs.~\eqref{eq:T}--\eqref{eq:magic} produces
all seven magic numbers exactly, with $K = 4$ as the unique
solution. The stability ceiling $K^K = 256$ correctly predicts
the endpoint of nuclear stability.
The spin-orbit inversion is a geometric theorem.

\paragraph{What is not established.}
The $Z$-dependent Coulomb correction (vortex repulsion among
proton fields) is not separated in the current formula.
Incorporating it would require a second term proportional to
$Z(Z{-}1)/A^{1/3}$, reducing the heavy-element overprediction.
However, the Coulomb interaction is not additive in this
framework: previous analysis shows that subtracting the
standard Coulomb term \emph{degrades} the fit, suggesting
the electromagnetic contribution is already partially encoded
in the constitutive law~\cite{nelson2025b}.

The element-specific cavity resonance structure (golden-ratio
harmonics within TetB) produces $\sim\!5\%$ element-by-element
variation in the effective $C$-factor. This is documented
but not included in the present letter.

\paragraph{Reproducibility.}
A self-contained Python validation script and an interactive
periodic table tool are available as supplementary material
at~\cite{zenodo2025}. The script requires no external libraries
and runs the complete 82-isotope validation in under one second.


% ================================================================
\section*{Acknowledgments}
% ================================================================
The author thanks the SPARC collaboration for public rotation
curve data and the AME2020 collaboration for nuclear mass
evaluations.


% ================================================================
\begin{thebibliography}{12}

\bibitem{nelson2025a}
S.~Nelson,
``Gravity Field Dynamics: Covariant completion of the
  scalar-tensor action from tetrahedral topology,''
  \textit{submitted to Class.\ Quantum Grav.} (2025);
  framework documentation at \texttt{docs.fielddynamics.org}.

\bibitem{nelson2025b}
S.~Nelson,
``GFD Nuclear Enterprise Validation,''
  Zenodo (2025), DOI: [to be assigned].

\bibitem{regge1961}
T.~Regge,
``General relativity without coordinates,''
  \textit{Nuovo Cimento} \textbf{19}, 558 (1961).

\bibitem{hassanrosen2012}
S.~F.~Hassan and R.~A.~Rosen,
``Bimetric gravity from ghost-free massive gravity,''
  \textit{J.\ High Energy Phys.} \textbf{2012}, 126 (2012).

\bibitem{boulwaredeser1972}
D.~G.~Boulware and S.~Deser,
``Can gravitation have a finite range?''
  \textit{Phys.\ Rev.\ D} \textbf{6}, 3368 (1972).

\bibitem{weinberg1995}
S.~Weinberg,
\textit{The Quantum Theory of Fields}, Vol.~1
  (Cambridge University Press, 1995).

\bibitem{ashtekar1986}
A.~Ashtekar,
``New variables for classical and quantum gravity,''
  \textit{Phys.\ Rev.\ Lett.} \textbf{57}, 2244 (1986).

\bibitem{bekenstein1984}
J.~D.~Bekenstein and M.~Milgrom,
``Does the missing mass problem signal the breakdown of
  Newtonian gravity?''
  \textit{Astrophys.\ J.} \textbf{286}, 7 (1984).

\bibitem{zenodo2025}
S.~Nelson,
``GFD Matter Profile: validation scripts and interactive tool,''
  Zenodo (2025), DOI: [to be assigned].

\bibitem{ame2020}
W.~J.~Huang \textit{et al.},
``The AME 2020 atomic mass evaluation,''
  \textit{Chinese Phys.\ C} \textbf{45}, 030002 (2021).

\bibitem{goeppert1949}
M.~Goeppert~Mayer,
``On closed shells in nuclei. II,''
  \textit{Phys.\ Rev.} \textbf{75}, 1969 (1949).

\end{thebibliography}


% ================================================================
% FULL TABLE (switch to one-column for wide table)
% ================================================================
\onecolumn

\begin{longtable}{r l r r r r r r r r r r}
\caption{GFD matter profile for all 83 stable elements. $A$ is the
most abundant stable isotope mass number. Shell~$n$ from
Eq.~\protect\eqref{eq:magic}. $\mu_2$ and $s_2$ from
Eq.~\protect\eqref{eq:mu2}. BE/A predicted from
Eq.~\protect\eqref{eq:binding} with $C = 2.172$~MeV. Observed
values from AME2020~\protect\cite{ame2020}.
$\Delta$ is the ceiling fraction $A/256$.}
\label{tab:full} \\
\toprule
$Z$ & Sym & $A$ & $n$ & $P(n)$ & $M(n)$ & $\mu_2$ & $s_2$ &
\multicolumn{2}{c}{BE/A (MeV)} & Err & $\Delta$ \\
& & & & & & & & GFD & Obs & (\%) & (\%) \\
\midrule
\endfirsthead
\toprule
$Z$ & Sym & $A$ & $n$ & $P(n)$ & $M(n)$ & $\mu_2$ & $s_2$ &
\multicolumn{2}{c}{BE/A (MeV)} & Err & $\Delta$ \\
& & & & & & & & GFD & Obs & (\%) & (\%) \\
\midrule
\endhead
\bottomrule
\endfoot
 1 & H  &   1 & 1 &   1 &   2 & 0.382 & $-$0.236 &  0.00 &  0.00 &     & 0.4 \\
 2 & He &   4 & 2 &   4 &   8 & 0.779 & $+$0.559 &  4.85 &  7.07 & $-$31 & 1.6 \\
 3 & Li &   7 & 2 &   4 &   8 & 0.867 & $+$0.734 &  6.38 &  5.61 & $+$14 & 2.7 \\
 4 & Be &   9 & 3 &  10 &  20 & 0.895 & $+$0.790 &  6.86 &  6.46 & $+$6  & 3.5 \\
 5 & B  &  11 & 3 &  10 &  20 & 0.913 & $+$0.826 &  7.18 &  6.93 & $+$4  & 4.3 \\
 6 & C  &  12 & 3 &  10 &  20 & 0.920 & $+$0.840 &  7.30 &  7.68 & $-$5  & 4.7 \\
 7 & N  &  14 & 3 &  10 &  20 & 0.932 & $+$0.865 &  7.49 &  7.48 & $+$0  & 5.5 \\
 8 & O  &  16 & 3 &  10 &  20 & 0.939 & $+$0.879 &  7.63 &  7.98 & $-$4  & 6.3 \\
 9 & F  &  19 & 3 &  10 &  20 & 0.949 & $+$0.898 &  7.80 &  7.78 & $+$0  & 7.4 \\
10 & Ne &  20 & 3 &  10 &  20 & 0.951 & $+$0.903 &  7.84 &  8.03 & $-$2  & 7.8 \\
11 & Na &  23 & 4 &  14 &  28 & 0.957 & $+$0.914 &  7.94 &  8.11 & $-$2  & 9.0 \\
12 & Mg &  24 & 4 &  14 &  28 & 0.959 & $+$0.918 &  7.97 &  8.26 & $-$4  & 9.4 \\
13 & Al &  27 & 4 &  14 &  28 & 0.964 & $+$0.927 &  8.05 &  8.33 & $-$3  & 10.5 \\
14 & Si &  28 & 4 &  14 &  28 & 0.965 & $+$0.930 &  8.08 &  8.45 & $-$4  & 10.9 \\
15 & P  &  31 & 5 &  25 &  50 & 0.968 & $+$0.936 &  8.13 &  8.48 & $-$4  & 12.1 \\
16 & S  &  32 & 5 &  25 &  50 & 0.969 & $+$0.938 &  8.15 &  8.49 & $-$4  & 12.5 \\
17 & Cl &  35 & 5 &  25 &  50 & 0.972 & $+$0.943 &  8.19 &  8.52 & $-$4  & 13.7 \\
18 & Ar &  40 & 5 &  25 &  50 & 0.975 & $+$0.951 &  8.26 &  8.60 & $-$4  & 15.6 \\
19 & K  &  39 & 5 &  25 &  50 & 0.975 & $+$0.949 &  8.24 &  8.56 & $-$4  & 15.2 \\
20 & Ca &  40 & 5 &  25 &  50 & 0.975 & $+$0.951 &  8.26 &  8.55 & $-$3  & 15.6 \\
21 & Sc &  45 & 5 &  25 &  50 & 0.978 & $+$0.956 &  8.30 &  8.62 & $-$4  & 17.6 \\
22 & Ti &  48 & 5 &  25 &  50 & 0.979 & $+$0.959 &  8.33 &  8.72 & $-$5  & 18.8 \\
23 & V  &  51 & 6 &  41 &  82 & 0.981 & $+$0.961 &  8.35 &  8.74 & $-$5  & 19.9 \\
24 & Cr &  52 & 6 &  41 &  82 & 0.981 & $+$0.962 &  8.35 &  8.78 & $-$5  & 20.3 \\
25 & Mn &  55 & 6 &  41 &  82 & 0.982 & $+$0.964 &  8.37 &  8.77 & $-$5  & 21.5 \\
26 & Fe &  56 & 6 &  41 &  82 & 0.982 & $+$0.965 &  8.38 &  8.79 & $-$5  & 21.9 \\
27 & Co &  59 & 6 &  41 &  82 & 0.983 & $+$0.966 &  8.39 &  8.77 & $-$4  & 23.0 \\
28 & Ni &  58 & 6 &  41 &  82 & 0.983 & $+$0.966 &  8.39 &  8.73 & $-$4  & 22.7 \\
29 & Cu &  63 & 6 &  41 &  82 & 0.984 & $+$0.969 &  8.41 &  8.75 & $-$4  & 24.6 \\
30 & Zn &  64 & 6 &  41 &  82 & 0.984 & $+$0.969 &  8.41 &  8.74 & $-$4  & 25.0 \\
31 & Ga &  69 & 6 &  41 &  82 & 0.986 & $+$0.971 &  8.43 &  8.72 & $-$3  & 27.0 \\
32 & Ge &  74 & 6 &  41 &  82 & 0.987 & $+$0.973 &  8.45 &  8.71 & $-$3  & 28.9 \\
33 & As &  75 & 6 &  41 &  82 & 0.987 & $+$0.974 &  8.45 &  8.70 & $-$3  & 29.3 \\
34 & Se &  80 & 6 &  41 &  82 & 0.988 & $+$0.975 &  8.47 &  8.71 & $-$3  & 31.3 \\
35 & Br &  79 & 6 &  41 &  82 & 0.987 & $+$0.975 &  8.47 &  8.70 & $-$3  & 30.9 \\
36 & Kr &  84 & 7 &  63 & 126 & 0.988 & $+$0.976 &  8.48 &  8.72 & $-$3  & 32.8 \\
37 & Rb &  85 & 7 &  63 & 126 & 0.988 & $+$0.977 &  8.48 &  8.70 & $-$3  & 33.2 \\
38 & Sr &  88 & 7 &  63 & 126 & 0.989 & $+$0.977 &  8.49 &  8.73 & $-$3  & 34.4 \\
39 & Y  &  89 & 7 &  63 & 126 & 0.989 & $+$0.978 &  8.49 &  8.71 & $-$3  & 34.8 \\
40 & Zr &  90 & 7 &  63 & 126 & 0.989 & $+$0.978 &  8.49 &  8.71 & $-$3  & 35.2 \\
41 & Nb &  93 & 7 &  63 & 126 & 0.989 & $+$0.979 &  8.50 &  8.66 & $-$2  & 36.3 \\
42 & Mo &  98 & 7 &  63 & 126 & 0.990 & $+$0.980 &  8.51 &  8.64 & $-$1  & 38.3 \\
43 & Tc &  98 & 7 &  63 & 126 & 0.990 & $+$0.980 &  8.51 &  8.59 & $-$1  & 38.3 \\
44 & Ru & 102 & 7 &  63 & 126 & 0.990 & $+$0.980 &  8.51 &  8.61 & $-$1  & 39.8 \\
45 & Rh & 103 & 7 &  63 & 126 & 0.990 & $+$0.981 &  8.52 &  8.58 & $-$1  & 40.2 \\
46 & Pd & 106 & 7 &  63 & 126 & 0.991 & $+$0.981 &  8.52 &  8.57 & $-$1  & 41.4 \\
47 & Ag & 107 & 7 &  63 & 126 & 0.991 & $+$0.981 &  8.52 &  8.55 &  $0$  & 41.8 \\
48 & Cd & 114 & 7 &  63 & 126 & 0.991 & $+$0.982 &  8.53 &  8.54 &  $0$  & 44.5 \\
49 & In & 115 & 7 &  63 & 126 & 0.991 & $+$0.983 &  8.53 &  8.52 &  $0$  & 44.9 \\
50 & Sn & 120 & 7 &  63 & 126 & 0.992 & $+$0.983 &  8.54 &  8.51 &  $0$  & 46.9 \\
51 & Sb & 121 & 7 &  63 & 126 & 0.992 & $+$0.984 &  8.54 &  8.48 & $+$1  & 47.3 \\
52 & Te & 130 & 8 &  92 & 184 & 0.992 & $+$0.985 &  8.55 &  8.47 & $+$1  & 50.8 \\
53 & I  & 127 & 8 &  92 & 184 & 0.992 & $+$0.984 &  8.55 &  8.45 & $+$1  & 49.6 \\
54 & Xe & 132 & 8 &  92 & 184 & 0.992 & $+$0.985 &  8.56 &  8.44 & $+$1  & 51.6 \\
55 & Cs & 133 & 8 &  92 & 184 & 0.993 & $+$0.985 &  8.56 &  8.41 & $+$2  & 52.0 \\
56 & Ba & 138 & 8 &  92 & 184 & 0.993 & $+$0.986 &  8.56 &  8.39 & $+$2  & 53.9 \\
57 & La & 139 & 8 &  92 & 184 & 0.993 & $+$0.986 &  8.56 &  8.38 & $+$2  & 54.3 \\
58 & Ce & 140 & 8 &  92 & 184 & 0.993 & $+$0.986 &  8.57 &  8.38 & $+$2  & 54.7 \\
59 & Pr & 141 & 8 &  92 & 184 & 0.993 & $+$0.986 &  8.57 &  8.35 & $+$3  & 55.1 \\
60 & Nd & 142 & 8 &  92 & 184 & 0.993 & $+$0.986 &  8.57 &  8.35 & $+$3  & 55.5 \\
61 & Pm & 145 & 8 &  92 & 184 & 0.993 & $+$0.986 &  8.57 &  8.32 & $+$3  & 56.6 \\
62 & Sm & 152 & 8 &  92 & 184 & 0.993 & $+$0.987 &  8.57 &  8.30 & $+$3  & 59.4 \\
63 & Eu & 153 & 8 &  92 & 184 & 0.993 & $+$0.987 &  8.57 &  8.28 & $+$4  & 59.8 \\
64 & Gd & 158 & 8 &  92 & 184 & 0.994 & $+$0.987 &  8.57 &  8.27 & $+$4  & 61.7 \\
65 & Tb & 159 & 8 &  92 & 184 & 0.994 & $+$0.987 &  8.57 &  8.25 & $+$4  & 62.1 \\
66 & Dy & 164 & 8 &  92 & 184 & 0.994 & $+$0.988 &  8.58 &  8.25 & $+$4  & 64.1 \\
67 & Ho & 165 & 8 &  92 & 184 & 0.994 & $+$0.988 &  8.58 &  8.23 & $+$4  & 64.5 \\
68 & Er & 166 & 8 &  92 & 184 & 0.994 & $+$0.988 &  8.58 &  8.23 & $+$4  & 64.8 \\
69 & Tm & 169 & 8 &  92 & 184 & 0.994 & $+$0.988 &  8.58 &  8.21 & $+$5  & 66.0 \\
70 & Yb & 174 & 8 &  92 & 184 & 0.994 & $+$0.989 &  8.58 &  8.20 & $+$5  & 68.0 \\
71 & Lu & 175 & 8 &  92 & 184 & 0.994 & $+$0.989 &  8.59 &  8.19 & $+$5  & 68.4 \\
72 & Hf & 180 & 8 &  92 & 184 & 0.994 & $+$0.989 &  8.59 &  8.18 & $+$5  & 70.3 \\
73 & Ta & 181 & 8 &  92 & 184 & 0.994 & $+$0.989 &  8.59 &  8.17 & $+$5  & 70.7 \\
74 & W  & 184 & 8 &  92 & 184 & 0.995 & $+$0.989 &  8.59 &  8.15 & $+$5  & 71.9 \\
75 & Re & 187 & 9 & 129 & 258 & 0.995 & $+$0.989 &  8.59 &  8.13 & $+$6  & 73.0 \\
76 & Os & 192 & 9 & 129 & 258 & 0.995 & $+$0.990 &  8.59 &  8.12 & $+$6  & 75.0 \\
77 & Ir & 193 & 9 & 129 & 258 & 0.995 & $+$0.990 &  8.59 &  8.11 & $+$6  & 75.4 \\
78 & Pt & 195 & 9 & 129 & 258 & 0.995 & $+$0.990 &  8.60 &  8.10 & $+$6  & 76.2 \\
79 & Au & 197 & 9 & 129 & 258 & 0.995 & $+$0.990 &  8.60 &  8.08 & $+$6  & 77.0 \\
80 & Hg & 202 & 9 & 129 & 258 & 0.995 & $+$0.990 &  8.60 &  8.06 & $+$7  & 78.9 \\
81 & Tl & 205 & 9 & 129 & 258 & 0.995 & $+$0.990 &  8.60 &  8.05 & $+$7  & 80.1 \\
82 & Pb & 208 & 9 & 129 & 258 & 0.995 & $+$0.990 &  8.60 &  7.87 & $+$9  & 81.3 \\
83 & Bi & 209 & 9 & 129 & 258 & 0.995 & $+$0.990 &  8.60 &  7.85 & $+$10 & 81.6 \\
\end{longtable}

\end{document}
